% ID: SEL-D0013-BC-E06
% ESTADO: cotejada. Segunda pasada visual completa y borrador compilado revisados.
% Fuente: Selectividad/matematicas_aplicadas/matematicas_aplicadas_P_A.pdf
% SHA-256: 0305750f6b6c3cc3edbba71a19cd7cce4276935dff5a0b7d964b02eaedeaa222
% Página: 2; bloque C; ejercicio original 6.

Una empresa fabrica dos tipos de bombillas: una LED y otra halógena. Se sabe
que un 5\% de las LED y un 2\% de las halógenas salen defectuosas. Se elige al
azar una bombilla de una caja que contiene 40 bombillas LED y 10 halógenas.
\begin{enumerate}[label=\alph*)]
  \item \textbf{(1.5 puntos)} Calcule la probabilidad de que la bombilla
  elegida no sea defectuosa.
  \item \textbf{(1 punto)} Calcule la probabilidad de que la bombilla elegida
  sea LED, sabiendo que es defectuosa.
\end{enumerate}
